\documentclass[a4paper, 11pt, oneside]{book}
\usepackage{graphicx}
\usepackage[utf8]{inputenc}
\usepackage[T1]{fontenc} 
\usepackage{fouriernc}
\usepackage[document]{ragged2e}
\usepackage{amsmath}
\usepackage{amssymb}
\graphicspath{ {images/} }
\usepackage{hyperref}
\usepackage[normalem]{ulem}
\usepackage[]{algorithm2e}
\usepackage{pgf,tikz}
\usetikzlibrary{spy}
\usetikzlibrary{backgrounds}
\usetikzlibrary{decorations.pathreplacing}
\usetikzlibrary{intersections}
\usepackage[margin=0.5in]{geometry}
\usepackage{textcomp}
\usetikzlibrary {positioning}
\usepackage{fancyhdr}
\pagestyle{fancy}
\fancyhead{}
\fancyhead[RO,LE]{Thesis Title}
\fancyfoot{}
\fancyfoot[LE,RO]{\thepage}
\fancyfoot[LO,CE]{Chapter \thechapter}
\fancyfoot[CO,RE]{Ritam Raha}

\usepackage{amsthm}
\theoremstyle{definition}
\newtheorem{theorem}{Theorem}[section]
\newtheorem{claim}[theorem]{Claim}
\newtheorem{conjecture}[theorem]{Conjecture}
\newtheorem{corollary}[theorem]{Corollary}
\newtheorem{definition}[theorem]{Definition}
\newtheorem{example}[theorem]{Example}
\newtheorem{fact}[theorem]{Fact}
\newtheorem{lemma}[theorem]{Lemma}
\newtheorem{proposition}[theorem]{Proposition}
\newtheorem{remark}[theorem]{Remark}




\begin{document}
\begin{titlepage}
	\centering
	\scshape 
	{\LARGE{CHENNAI MATHEMATICAL INSTITUTE}}\\
          \vskip 0.5cm
	MASTERS THESIS
	\vspace*{\baselineskip} 
	%------------------------------------------------
	%	Title
	%------------------------------------------------
	
	\rule{\textwidth}{1.6pt}\vspace*{-\baselineskip}\vspace*{2pt}
	\rule{\textwidth}{0.4pt} 

	\vspace{0.75\baselineskip}
	
	{\LARGE THE NAME\\ that is \\ YET TO BE DECIDED\\} 
	
	\vspace{0.75\baselineskip} 
	
	\rule{\textwidth}{0.4pt}\vspace*{-\baselineskip}\vspace{3.2pt} 
	\rule{\textwidth}{1.6pt}
	
	\vspace{2\baselineskip} 
	
	%------------------------------------------------
	%	Author
	%------------------------------------------------
	
	
	
	\textit{Author:}
	
	\vspace{0.5\baselineskip} 
	
	{\scshape\Large Ritam Raha}\\ % Author
	\vspace{0.2\baselineskip}
	{\small Chennai Mathematical Institute}
	\vspace{2\baselineskip} 

		
	\textit{Supervised by:}     
	\vspace{0.5\baselineskip}

	{\scshape\Large Nicolas Markey}\ \ \ \ \&\ \ \ \  {\scshape\Large Loïc Hélouët}\\ % Supervisors
	\vspace{0.2\baselineskip}
	{\small INRIA Rennes}
	\vspace{2\baselineskip}

	\vspace{2\baselineskip} 
	
	\textit{A thesis submitted in fulfillment of the requirements\\ for the degree of Master of Science}\\
	\vspace{1\baselineskip}
	\textit{in the}\\
	\vspace{1\baselineskip}
	{\scshape Department of Computer Science\\ Chennai Mathematical Institute}

	\vspace{4\baselineskip}
	
	\includegraphics[width=0.7\textwidth]{cmi.png}

	\vfill % Whitespace between editor names and publisher logo
	
\end{titlepage}


\chapter*{Abstract}
Quantitative reachability games are finite two player turn-based games played on weighted graphs. The objective of the game combines reachability objective(qualitative) with the (quantitative) requirement that the weights along a path must satisfy certain constraints (bounds). Beside having direct applications in reactive system synthesis with resource constraints, it is one of the simplest models that combine quantitative and qualitative objectives. In this thesis, we ask the question whether Player $1$ can  reach a target state along a path for which the accumulated weight for any finite prefix satisfies some bound. We look both at single bound and dual bound problems. We also consider the weighted game where we weaken one of the bounds (e.g. we allow violations of a given bound but up to a bounded number of times).
 
\chapter*{Dedication}
To 
 
\chapter*{Declaration}
I declare that..
 
\chapter*{Acknowledgements}
I want to thank...

\tableofcontents

\chapter{Introduction}
\paragraph{Games on weighted graphs.}
Weighted games are a common way to formally address questions
related to consumption, production and storage of resources: the arena
of such game are two-player turn-based games in which transitions
carry positive or negative integers, representing the accumulation or
consumption of resource.  Various objectives have been considered for
such arenas, such as optimizing the total or average amount of
resources that have been collected along the play, or maintaining the
total amount within given bounds. The~latter kind of objectives,
usually referred to as \emph{energy
objectives}~\cite{CdAHS03,BouyerFLMS08}, has been widely studied in
the untimed
setting~\cite{ChatterjeeD12,ChatterjeeDHR10,DDGRT10,FahrenbergJLS11,JLR13,
JLS15,VCDHRR15,BMRLL15,BHMRZ17,DM18}, and to a lesser extent in the
timed setting~\cite{BFLM10,qest2012-BLM}.
%
As their name indicates, energy objectives can be used to model the
evolution of the available energy in an autonomous system: besides
achieving its tasks, the~system has to take care of recharging
batteries regularly enough so as to never run out of power.  Energy
objectives were also used to model moulding machines: such machines
inject molten plastic into a mould, using pressure obtained by storing
liquid in a tank~\cite{CJLRR09}; the~level of liquid has to be
controlled in such a way that enough pressure is always available, but
excessive pressure in the tank would reduce the service life of the valve.

Energy games impose strict constraints on the total amount of energy
at all stages of the play. Two kinds of constraints have been mainly
considered in the literature: lower-bound constraints (a.k.a. L-energy constraints) impose a strict lower bound (usually~zero), but impose no
upper bound; on~the other hand, lower- and upper-bound constraints
(a.k.a. LU-energy constraints) require that the energy level always
remain within a bounded interval~$[L;U]$. Finding strategies that
realize L-energy objectives along infinite runs is in PTIME in the
one-player setting, and in NP $\cap$ coNP for two players;
for LU-energy objectives, it~is respectively PSPACE-complete
and EXPTIME-complete~\cite{BouyerFLMS08}. Some works have also considered
the existence of an initial energy level for which a winning strategy
exist~\cite{ChatterjeeDHR10}.
%
Energy objectives have also been combined with other objectives,
either qualitative (e.g.~parity~\cite{ChatterjeeD12}) or quantitative
objectives (e.g.~multi-dimensional
energy~\cite{ChatterjeeDHR10,FahrenbergJLS11,JLR13}). 

In this paper, we focus on weighted games combining energy objectives
together with reachability objectives. Our~first result is the
(expected) proof that L-energy games with or without reachability
objectives are interreducible; the same holds for LU-energy games.

%
We~then focus on relaxations of the energy constraints, in two
different directions. In both cases, the lower bound remains
unchanged, as it corresponds to running out of energy, which we always
want to avoid. We~thus only relax the upper-bound constraint.
The~first direction concerns \emph{weak upper bounds}, already
introduced in~\cite{BouyerFLMS08}: in~that setting, hitting the upper bound
is allowed, but there is no overload, i.e. trying to exceed the upper bound will simply maintain the energy level at this maximal level. Yet, a strict lower bound is
still imposed. We~name these objectives LW-energy
objectives. This~could be used as a (simplified) model for batteries.
When considered alone, LW-energy objectives are not much different
from~ L-energy objectives, in the sense that the aim is to find a
reachable \emph{positive loop}. LW-energy games are in PTIME for
one-player games, and in NP $\cap$ coNP for two players. When
combining LW-energy and reachability objectives, the~situation
changes: different loops may have different effects on the energy
level, and we have to keep track of the final energy level reached
when iterating those loops.

The second way of relaxing upper bounds, which we call \emph{soft
upper bound}, allows a limited number (or amount) of
violations: when modeling a pressure tank, the lower-bound constraint
is strict (pressure should always be available) but the upper bound is
soft (excessive pressure may be temporarily allowed is
needed). We~consider different kinds of limits (on the number or
amount of violations), and prove that (with or without reachability
objectives) the energy game problems are PSPACE-complete for one-player arenas,
and EXPTIME-complete for two-player ones.

\fbox{minimization?}


\paragraph{Related work.}
Quantitative games have been the focus of numerous research articles
 since the~1970s, with various kinds of objectives. 
%such as ultimately
% optimizing the total payoff~\cite{}, mean-payoff~\cite{EM79,ZwickP95},
% or discounted sum~\cite{ZwickP95,Andersson06}.
% %
% Energy objectives, which are a kind of safety objectives on the total
% payoff, were introduced in~\cite{CdAHS03} and rediscovered
% in~\cite{BouyerFLMS08}.
% %
% Several works have extended those works by
% combining quantitative conditions together, e.g. multi-dimensional
% energy conditions~\cite{FahrenbergJLS11,JurdzinskiLS15} or
% conjunctions of energy- and mean-payoff
% objectives~\cite{ChatterjeeDHR10}. Combinations with qualitative
% objectives (e.g. reachability~\cite{Chatterjee0H17} or parity
% objectives~\cite{ChatterjeeD12,ChatterjeeRR14}) were also considered.
% %
% Similar objectives have been considered in slightly different
% settings e.g. Vector Addition Systems with States~\cite{Reichert16} and
% one-counter machines~\cite{GHOW10,Hun15}.
% %



Other types of quantitative games consider mean payoffs. The payoff
of a run is the mean weight along this run, and the value of a
mean payoff game is the maximal/minimal mean weight that one can
enforce with an appropriate strategy. In this setting, one mainly has
to consider the mean value of cycles, that absorb all other
costs. \cite{ZwickP95} shows that the value of infinite plays can be
computed in $??$, where $p$ is the
maximal absolute value of weights in the arena.

Mean payoff games, however, only address quantitative questions as
limit values of infinite runs, and do not consider quantitative
constraints on prefixes of these runs, nor boolean objectives.

Discounted games are another form of quantitative games for which
polynomial solutions exist. In these games, a discounting
factor~$\lambda < 1$ is defined, and the contribution of the $i$-th
transition to the total payoff of a run is the weight of the
transition multiplied by $\lambda^i$. The value of these games can be
computed in PTIME using linear programming~\cite{Andersson06}.



Games with quantitative and reachability objectives have already been
considered. \cite{Chatterjee0H17} considers total payoff, discounted
payoff and energy payoff combined with a reachability objective. Given
a game graph $G$ an initial state~$s$, and a target state~$t$, the
question addressed is whether there exists a~run~$\rho$ from $s$
to~$t$ in $G$ such that the total payoff, discounted payoff, or the
energy level of $\rho$ is non negative. Checking that a run with
payoff $\geq 0$ exists for a single player, is~in~ PTIME for all payoff functions, and for
two-player games in NP $\cap$ coNP. For equality, the problems are
respectively NP-complete and EXPSPACE-complete for total and
energy payoff, and the question is still open for discounted
games. The~PTIME algorithms rely on the computation of values in
games, which can be efficiently done as soon as one does not impose
constraints on payoff achieved by prefixes of winning runs.

\cite{CdAHS03} combines energy objectives and B\"uchi objectives in
two players games; the objective is that a pair of components sharing
resources never consume more than a certain threshold, while
preserving some liveness properties, i.e. visiting some states
infinitely often.

\cite{ChatterjeeD12} combines quantitative and boolean objectives in
Energy-Parity Games: the objective is to win a parity game while
maintaining an energy constraint, i.e. ensure that the total payoff of
every prefix of a winning run remains $\geq 0$. Results are the following: first player 1 has winning strategies with memory of
size $n\cdot d \cdot W$, where $n$ is the number of states, $W$ the
largest weight, and $d$ the number of priorities in the parity
condition. Second, memoryless strategies are sufficient for player 2. Last
memoryless winning strategies exist for player 2 if the boolean condition
is a coB\"uchi condition.

\cite{ChatterjeeDHR10} generalizes mean payoff and energy games, by considering vectors of payoff functions.
These games are determined when considering finite-memory
strategies. In this setting, energy and mean payoff games are
interreducible, and threshold questions for mean payoff, or unknown
initial credit (whether there exist initial values for which an energy
game is winning) are coNP-complete.

\cite{ChatterjeeRR14} considers multiple quantitative mean payoff and
energy objectives
combined with parity objectives.  Strategies for such quantitative
games have to fulfill two objectives: first progress towards a target
state or enforce a parity condition, and second, satisfy the
constraint on a vector of quantitative outcomes of the runs: all mean
payoff are greater than a threshold $\nu$, or all energy levels are
positive. Usually these games require infinite memory. However, when
restricting to finite memory strategies, the required memory is
exponential.

\cite{BouyerFLMS08} considers energy games with strong lower bound
constraints (L), strong lower bound and strong upper bound (LU) and strong
lower bound and weak upper bound (LW) constraints. Winning runs in these
games are infinite runs which prefixes satisfy the L, LU or LW energy
constraint. For the one player setting, existence of strategies for L,
and LW games are in PTIME, and LU games are PSPACE-complete. In the two-player setting L, and LW games are in
NP $\cap$ coNP, and LU games are EXPTIME-complete.

\cite{FahrenbergJLS11} considers energy games in multi-weighted automata,
and L, LW, LU games similar to those of \cite{BouyerFLMS08} but with
vectors of energy levels and constraints attached to each energy level
in these vectors. The complexity of these games depend on the size of
energy vectors.

\cite{HunterR14} considers multiple quantitative objectives,
i.e. addresses the problem of determining if a player can attain a
payoff in a finite union of arbitrary intervals for various payoff
functions (liminf, mean-payoff, discount sum, total~sum). Given a
fixed number of intervals $I_1, \dots I_k$ and a payoff function~$f$,
a play~$\pi$ is winning if $f(\pi) \in I_j$ for some $j\in [1;k]$. The
complexity of these games depend on the payoff function and ranges from
NP $\cap$ coNP to EXPSPACE.


\cite{JurdzinskiLS15} considers games in multi-weighted automata and
the initial credit problem, i.e. whether, starting from a vertex~$v$,
there exists a vector~$B\in \mathbb N^k$ of initial weights such that player 1 has a
winning strategy starting from configuration $(v,B)$. They show that
this problem is 2EXPTIME-complete for general multi-dimensional energy
games.

\cite{Reichert16} Considers reachability in counter games under different semantics. Configurations of the game are counter values, and moves are labeled by integer vectors. The semantics of the arena either imposes no constraint on legal values of counters ($\mathbb{ Z}$ semantics), forbids moves that would result in a negative value of some counter (blocking VASS semantics), or considers VASS semantics with weak lower bounds (all moves are legal, but counter values have $0$ lower bound). A play is winning if it reaches a configuration from a target set. In dimension 2, under every semantics, two-player reachability games are undecidable. In dimension 1, these games are in EXPSPACE. 




 
\chapter{Preliminaries}
\section{Definitions}

\textbf{Game graphs.}  A \textit{game graph} $G= \langle Q,E,w \rangle$ consists of a finite set $Q$ of states partitioned in to player-$1$ states $Q_1$ and player-$2$ states $Q_2$ (i.e., $Q = Q_1 \cup Q_2$), and a set $E \subseteq Q \times Q$ of edges such that for all $q \in Q$, there exists (at least one) $q^{\prime} \in Q$ such that $(q, q^{\prime}) \in E$. A player-$1$ game is a game graph where $Q_1 = Q$ and $Q_2 = \phi$. A special subset $T \subseteq Q$ of vertices is known as the \textit{Target states} or the \textit{Goal states}. Now, $w : E \rightarrow \mathbb{R}$ is the weight function such that $w(q_j,q_{j+1})$ is the weight of the edge between the vertices $q_j$ and $q_{j+1}$ \\
\vskip 0.6cm


\textbf{Plays and strategies.} A game on $G$ starting from a state $q_0 \in Q$ is played in rounds as follows. If the game is in a player-$1$ state, then player $1$ chooses the successor state from the set of outgoing edges; otherwise the game is in a player-$2$ state, and player $2$ chooses the successor state. We always consider the reachability game, i.e. this ends as soon as any vertex from the set $T$ has been reached. The game results in a play from $q_0$ , i.e., a finite path $\rho = q_0 q_1 \ldots q_n$ such that $q_n \in T$ \& $(q_i , q_{i+1}) \in E$, $\forall i \geq 0$. \\
A strategy for player $1$ is a function $\sigma : Q^{*} Q_1 \rightarrow Q$ such that $(q, \sigma(\rho . q)) \in E$,
$\forall q \in Q_1$ and all $\rho \in Q^{*}$. An outcome of $\sigma$ from $q_0$ is a play $q_0 q_1 \ldots$ such that $\sigma(q_0 \ldots q_i) = q_ {i+1}$ for all $i \geq 0$ such that $q_i \in Q_1$ . Strategy and outcome for player 2 are defined analogously.\\
\vskip 0.6cm 


\textbf{Payoff functions.} We consider two different payoff functions:
\begin{itemize}
	\item \textit{Finite Total payoff.} For a finite path $p=q_0 q_1 \ldots q_l$, the total payoff of the path $p$ is defined as, $TP(p)= \sum_{i=0}^{l-1} w(q_i,q_{i+1})$.

	\item \textit{Finite Mean payoff.} For a finite path $p=q_0 q_1 \ldots q_l$, the mean payoff of the path $p$ is defined as, $MP(p)= (1/l) \cdot \sum_{i=0}^{l-1} w(q_i,q_{i+1})$.
\end{itemize}
\vskip 0.6cm


\textbf{Bounds on weights.} We consider two kinds of bounds, strong and week bounds:
\begin{itemize}
	\item \textit{Strong bounds.} This is the usual notion of bounds used in literature. Weight of a path $p$ is strongly\\
	 bounded(upper or lower) by $B$ means, for every finite prefix $\pi$ of the path $p$, $w(\pi) \gtreqless B$.

	\item \textit{Weak bounds.} This is a new notion of boundedness. Weight of a path $p= s_1 \rightarrow s_2 \rightarrow \ldots \rightarrow s_n$, starting from $s_1$ with $c \in \mathbb{R} \cup \{- \infty \}$ as a lower weak bound is defined inductively by $r_1 = max(0, c),\  r_{i+1} = max(r_i + w(s_i, s_{i+1}), c)$. The notion of weak upper bound is analogously defined.\\

	So for computing $w\downarrow_{c}(p)$, costs are accumulated along the transitions of $p$, but if at some point it goes down below $c$. it is reset to $c$ i.e. all possible decreases below $c$ are simply discarded.
\end{itemize}
\vskip 0.6cm

\textbf{Finite Memory\& Memoryless strategies.} A strategy for $P1$, $\sigma : Q^{*} Q_1 \rightarrow Q$ is called a \textit{finite-memory} strategy if every move depends on finite amount of history. The strategy is called a \textit{memoryless} one, if it does not depend on the whole history and only depends on the current state he is in. Hence, a memoryless strategy can be seen as a function $\sigma:Q_1 \rightarrow Q$. The definitions are analogous for $P2$.
\vskip 0.6cm

\textbf{Objectives.} In this thesis, we will focus on \textit{quantitative-reachability} objectives. We consider two kinds of quantitative functions $f: \varrho \rightarrow \mathbb{N}$, where $\varrho$ denotes the set of all finite paths in the corresponding game graph $G$\\
\begin{itemize}
\item[--] \textit{Total Payoff.} Total payoff function of a path $\rho= s_1 \rightarrow s_2 \rightarrow \ldots \rightarrow s_n$ is defined as $TP(\rho)= \sum_{i=1}^{n-1} w(s_i,s_i+1)$.\\
\item[--] \textit{Mean Payoff.} Mean payoff function of a path $\rho= s_1 \rightarrow s_2 \rightarrow \ldots \rightarrow s_n$ is defined as $MP(\rho)= (1/n)\sum_{i=1}^{n-1} w(s_i,s_i+1)$.\\

\end{itemize}
 We will look at the two kind of objectives:\\
 \begin{itemize}
 \item \textit{Single Bound Objectives.} Given a game graph $G$, a starting vertex $q_0$ and a target vertex $t$, a quantitative function $f$ and a bound $b \in \mathbb{N}$, single bound objective of $P1$ asks that, if starting from $q_0$, $P1$ can reach $t$ in a path $\rho$, such that $f(\rho) \lessgtr b$.
 
 \item \textit{Dual Bound Objectives.} Given a game graph $G$, a starting vertex $q_0$ and a target vertex $t$, a quantitative function $f$ and two bounds $U\  \& \ L \in \mathbb{N}$, single bound objective of $P1$ asks that, if starting from $q_0$, $P1$ can reach $t$ in a path $\rho$, such that $L \leq f(\rho) \leq U$. \\
\vskip 1cm
Note that, these bounds are kind of strong on the both upper and lower side. Hence, we are going to call these \textit{Strong Dual Bound Objectives}. Consequently, it is obvious that, we are going to define the notion of \textit{Weak Dual Bound Objectives} as follows:
\vskip 0.5cm
\textit{Weak Dual Bound Objectives.} Let $A$ be a weighted graph. $c \in \mathbb{R} \cup \{- \infty\}$
and let $\gamma =s_1 \rightarrow s_2 \rightarrow \cdots \rightarrow s_n \in Runs(A)$, starting from $s_1$. The accumulated weight with
initial credit 0 under weak lower bound $c$ is $w\mathord{\downarrow}_{c}(\gamma) = r_n$, where $r_1, \ldots, r_n \in \mathbb{R}$
are defined inductively by $r1 = max(0,c), r_{i+1} = max(r_i + w(s_i, s_{i+1}), c)$.\\
So for computing $w\downarrow_{c}(\gamma)$, costs are accumulated along the transitions of $\gamma$,
but if at some point it goes down below $c$. it is reset to $c$  i.e. all possible decreases below $c$ are
simply discarded. \\
 \end{itemize}
 \vskip 0.2cm
 
 \section{Thesis Description}
 In this thesis, we will look at the following games:\\
 In \textbf{Chapter Three}, we will see \textit{Quantitative-Reachability Games with Single Bound}, in \textbf{Chapter Four}, we will see \textit{Quantitative-Reachability Games with Dual Bounds}, both strong and weak objectives. In \textbf{Chapter Five}, we will see a new kind of game, \textit{Quantitative-reachability game with bounded number of violations} which we call as \textit{Apna Game}.
 
 

	

 
\chapter{Quantitative Reachability Games with Single Bounds}
In this chapter, we will see different quantitative games with single bound. \\

\section{Energy-Reachability Game}
Given a game graph $G=\langle Q_1, Q_2, E, w, q_0, T \rangle$, with all usual notations, the energy-reachability objective is to reach the target states $T$ starting from $q_0$ in a path $\rho$, such that for all finite prefixes $\pi(\rho)$, $w(\pi) \geq 0$.\\

Now, given a game graph, the decision problem for this game is that if Player $1$ can win maintaining energy-reachability objective? We will prove the following theorem:\\
\begin{theorem}
\label{energy-reach-thm}
The decision problem for Energy-Reachability game is in $NP \cap coNP$. For the lower bound, the problem is mean-payoff hard.
\end{theorem}
\begin{proof}
We will prove that, energy-reachability game is equivalent to infinite-energy game by both side reduction. The theorem will be the consequence of the proof from ~\cite{DBLP:conf/formats/BouyerFLMS08}.\\

Given a game graph $G=\langle Q_1, Q_2, E, w, q_0\rangle$, with all usual notations, the decision problem for infinite-energy game asks, if player $1$ has a strategy such that with zero initial weight for any infinite play $\gamma$ from initial state $q_0$, if $w(\gamma^{\prime}) \geq 0$ for all finite prefixes $\gamma^{\prime}$ of $\gamma$? \\
First we reduce the infinite-energy game to energy-reachability game.
Consider a game graph $G=\langle Q_1, Q_2, E, w, q_0\rangle$. We produce another game graph $G^{\prime}=\langle Q_1, Q_2, E^{\prime}, w^{\prime}, q_0,T\rangle$, where we add a target vertex $T$. For every player $1$ vertex $v \in Q_1$, we add an edge $(v,T) \in E^{\prime}$ with $w^{\prime}(v,T) = -\delta$, where $\delta$ is larger than the sum of all the positive weights of the graph $G$. For, all other edges $e \in E$, we have the same corresponding edges $e^{\prime} \in E^{\prime}$ with $w^{\prime}(e^{\prime})= w(e)+ \epsilon$, where $\epsilon << 1/2n$, where $n$ is the total number of vertices. The overview of the construction is shown in Figure \ref{inf-reach}.
\begin{figure}
    \label{inf-reach}
    \centering
    \tikzset{every picture/.style={line width=0.75pt}} %set default line width to 0.75pt        

\begin{tikzpicture}[x=0.75pt,y=0.75pt,yscale=-1,xscale=1]
%uncomment if require: \path (0,300); %set diagram left start at 0, and has height of 300

%Shape: Ellipse [id:dp35665433870720253] 
\draw   (28,115) .. controls (28,81.31) and (74.23,54) .. (131.25,54) .. controls (188.27,54) and (234.5,81.31) .. (234.5,115) .. controls (234.5,148.69) and (188.27,176) .. (131.25,176) .. controls (74.23,176) and (28,148.69) .. (28,115) -- cycle ;
%Shape: Circle [id:dp8502084996926933] 
\draw   (59,96.5) .. controls (59,89.04) and (65.04,83) .. (72.5,83) .. controls (79.96,83) and (86,89.04) .. (86,96.5) .. controls (86,103.96) and (79.96,110) .. (72.5,110) .. controls (65.04,110) and (59,103.96) .. (59,96.5) -- cycle ;
%Shape: Square [id:dp9450962306071624] 
\draw   (137,82) -- (159,82) -- (159,104) -- (137,104) -- cycle ;
%Straight Lines [id:da2842158542061044] 
\draw    (86,96.5) -- (135.5,93.14) ;
\draw [shift={(137.5,93)}, rotate = 536.11] [color={rgb, 255:red, 0; green, 0; blue, 0 }  ][line width=0.75]    (10.93,-3.29) .. controls (6.95,-1.4) and (3.31,-0.3) .. (0,0) .. controls (3.31,0.3) and (6.95,1.4) .. (10.93,3.29)   ;

%Shape: Ellipse [id:dp9633483807744296] 
\draw   (338,110) .. controls (334.57,76.31) and (378.02,49) .. (435.04,49) .. controls (492.07,49) and (541.08,76.31) .. (544.5,110) .. controls (547.93,143.69) and (504.48,171) .. (447.46,171) .. controls (390.43,171) and (341.42,143.69) .. (338,110) -- cycle ;
%Shape: Circle [id:dp3748033424118844] 
\draw   (367.12,91.5) .. controls (366.36,84.04) and (371.79,78) .. (379.24,78) .. controls (386.7,78) and (393.36,84.04) .. (394.12,91.5) .. controls (394.88,98.96) and (389.45,105) .. (381.99,105) .. controls (374.53,105) and (367.87,98.96) .. (367.12,91.5) -- cycle ;
%Shape: Square [id:dp9047184358152496] 
\draw   (443.64,77) -- (465.64,77) -- (467.88,99) -- (445.88,99) -- cycle ;
%Straight Lines [id:da04429014923239061] 
\draw    (394.12,91.5) -- (443.26,88.14) ;
\draw [shift={(445.26,88)}, rotate = 536.0899999999999] [color={rgb, 255:red, 0; green, 0; blue, 0 }  ][line width=0.75]    (10.93,-3.29) .. controls (6.95,-1.4) and (3.31,-0.3) .. (0,0) .. controls (3.31,0.3) and (6.95,1.4) .. (10.93,3.29)   ;

%Shape: Circle [id:dp6442851353260122] 
\draw   (428,227) .. controls (428,213.19) and (439.19,202) .. (453,202) .. controls (466.81,202) and (478,213.19) .. (478,227) .. controls (478,240.81) and (466.81,252) .. (453,252) .. controls (439.19,252) and (428,240.81) .. (428,227) -- cycle ;
%Straight Lines [id:da29377848783032556] 
\draw    (381.99,105) -- (440.49,205.27) ;
\draw [shift={(441.5,207)}, rotate = 239.74] [color={rgb, 255:red, 0; green, 0; blue, 0 }  ][line width=0.75]    (10.93,-3.29) .. controls (6.95,-1.4) and (3.31,-0.3) .. (0,0) .. controls (3.31,0.3) and (6.95,1.4) .. (10.93,3.29)   ;


% Text Node
\draw (107,83) node  [align=left] {w};
% Text Node
\draw (413.74,78) node  [align=left] {w + $\displaystyle \epsilon $};
% Text Node
\draw (453,227) node  [align=left] {T};
% Text Node
\draw (397,179) node  [align=left] {
\begin{itemize}
\item \mbox{-}$\displaystyle \delta $
\end{itemize}};
% Text Node
\draw (117,27) node [xslant=0.06] [align=left] {G};
% Text Node
\draw (436,18) node  [align=left] {G$\displaystyle ^{\prime }$};


\end{tikzpicture}

    \caption{Infinite-energy to Energy-reachability}
\end{figure}

If player $1$ has a strategy to win infinite-energy in $G$, then all the cycles player $1$ forces in the infinite winning path $\rho$ is either a positive or a zero cycle. The same path in $G^{\prime}$ will produce positive cycles for $\epsilon$. Hence, after a large number of iterations, player $1$, will have weight $\geq \delta$ in some of its vertex and then it can reach $T$ in $G^{\prime}$ and win energy-reachability. For the other side, let player $1$ has a winning strategy in energy-reachability in $G^{\prime}$. That means, in every winning path, at some player $1$ vertex it has weight $\geq \delta$, which is impossible without forcing a positive cycle as $\delta$ is large enough. Hence, iterating the same cycle player $1$ wins infinite-energy game in $G$.\\
\\
Now, we will show the opposite reduction i.e. from energy-reachability to infinite-energy. Again consider, a game graph $G=\langle Q_1, Q_2, E, w, q_0, T\rangle$. We produce another game graph $G^{\prime}=\langle Q_1, Q_2, E^{\prime}, w^{\prime}, q_0\rangle$, where we add a self loop of weight $0$ on $T$.\\
\huge\fcolorbox{blue}{red}{Incomplete}
\end{proof}

\section{Finite Mean-Payoff Reachability}
In this section, we will take the mean-payoff function as our quantitative function. In general mean-payoff is used in theory mostly for the case of infinite games, but in this thesis we restrict ourselves to reachability, hence finite games. For this reason, we will define a notion of mean payoff for the finite case:\\
\vskip 0.5cm
Given a game graph $G=\langle Q_1, Q_2, E, w, q_0,T\rangle$ and a bound $b \in \mathbb{N}$, the decision problem for finite mean-payoff reachability game asks, starting from $q_0$ with zero initial weight, if player $1$ has a strategy to reach $T$ in some path $\gamma$ such that, for all finite prefixes $\gamma^{\prime}$ of $\gamma$, if $MP(\gamma^{\prime}) \lesseqgtr b$? We will only consider the case of  $MP(\gamma^{\prime}) \leq b$ here. The other case is just the dual. Remember that, for a prefix $\gamma^{\prime}= q_0,q_1,\ldots q_l$, $MP(\gamma^{\prime})= (1/l) \cdot \Sigma_{i=0}^{l-1} w(q_i,q_i+1)$.\\ Now, we will prove the following result about the finite mean-payoff reachability game:\\
\\
\begin{theorem}
\label{fin-meanpayoff-thm}
The decision problem for finite mean-payoff reachability game is in $NP \cap coNP$ and also mean-payoff hard. 
\end{theorem}
\begin{proof}
We will reduce this game to and from energy-reachability game. As usual, we consider a game graph\\ $G=\langle Q_1, Q_2, E, w, q_0,T\rangle$. We produce another game graph, $G^{\prime}=\langle Q_1, Q_2, E, w^{\prime}= b-w (\text{or } w-b, \text{accordingly}), q_0,T\rangle$. We claim that, player $1$ wins energy-reachability in $G$ iff he wins finite mean-payoff reachability in $G^{\prime}$ and also vice-versa. This will prove that, energy-reachability is equivalent to mean-payoff reachability complexity-wise. The result of the theorem will be then the consequence of Theorem \ref{energy-reach-thm}\\
\\First we prove the reduction from energy reachability to mean-payoff reachability. The opposite reduction will be exactly similar. Let player $1$ has a winning strategy $\lambda$ in $G$ for energy reachability. That means, for all the path $\rho$ conforming $\lambda$, reaches $T$ and for all finite prefixes $\rho^{\prime}$ of $\rho$, $w(\rho)\geq 0$. Consider such a prefix $q_0,q_1,\ldots q_l$. We have, $\Sigma_{i=0}^{l-1} w(q_i,q_{i+1}) \geq 0$. Now, let in $G^{\prime}$, also player $1$ plays with same strategy $\lambda$. Now, for the same prefix, mean-payoff will be:
\begin{equation}
\begin{align*}
\notag
\centering
MP(\rho^{\prime})&= (1/l)\Sigma_{i=0}^{l-1} w^{\prime}(q_i,q_{i+1})\\
&=(1/l)\Sigma_{i=0}^{l-1} (b- w(q_i,q_{i+1}))\\
&=(1/l)(lb- \Sigma_{i=0}^{l-1} w(q_i,q_{i+1})\\
&=b- (1/l)\Sigma_{i=0}^{l-1} w(q_i,q_{i+1})\\
&\leq b
\end{align*}
\end{equation}

Hence, player $1$ wins in finite mean-payoff reachability objective in $G^{\prime}$. The other proofs can be obtained similarly.\\
Hence, we prove that, finite mean-payoff reachability is equivalent to energy-reachability.

\end{proof}

\section{Finite Total Payoff Reachability}
Now, it comes down to the final quantitative function we are concerned with in this thesis, which is total payoff. Given a game graph $G=\langle Q_1, Q_2, E, w, q_0,T\rangle$ and a bound $b \in \mathbb{N}$, the decision problem for finite total payoff reachability game asks, starting from $q_0$ with zero initial weight, if player $1$ has a strategy to reach $T$ in some path $\gamma$ such that, for all finite prefixes $\gamma^{\prime}$ of $\gamma$, if $TP(\gamma^{\prime}) \lesseqgtr b$? We will only consider the case of  $TP(\gamma^{\prime}) \geq b$ here. The other case is just the dual. Remember that, for a prefix $\gamma^{\prime}= q_0,q_1,\ldots q_l$, $TP(\gamma^{\prime})= \Sigma_{i=0}^{l-1} w(q_i,q_i+1)$. Note that, energy-reachability is just a special case of total payoff reachability, where $b=0$. We will show that, it is equivalent to energy-reachability game. From that, the following theorem is evident:\\
\begin{theorem}
\label{fin-totalpayoff-thm}
The decision problem for finite mean-payoff reachability game is in $NP \cap coNP$ and also mean-payoff hard. 
\end{theorem}
\begin{proof}
\huge\fcolorbox{blue}{red}{Incomplete}
\end{proof}

\section{Conclusion}
This brings us to the end of the Quantitative Reachability Games with Single Bounds. We showed that, all the reachability games with single bounds (for the quantitative functions we are concerned with in this thesis) are equivalent to energy-reachability game which is complexity-wise equivalent to infinite energy games. So, we have $NP \cap coNP$ complexity for all of them. For the lower bound also, all of them are at least mean-payoff hard. Next we move to the chapters for the similar kind of games with two bounds. 


 
\chapter{Chapter Four Title}
In this chapter, we will consider the quantitative games, where player $1$ has to reach his goal, always maintaining his weights inside two bounds. This chapter we will consider both the bounds to be strong.

\section{Finite Total Payoff Reachability}

In this section, we will use total payoff as our quantitative function. Given a game graph $G=\langle Q_1, Q_2, E, w, q_0, T \rangle$, an upper bound $U \in \mathbb{N}$ and a lower bound $b \in \mathbb{N}$, finite total payoff reachability objective with dual bound says that, with zero initial energy starting from $q_0$, player $1$ has to reach $T$ in a path $\rho$ such that, for all finite prefixes $\rho^{\prime}$ of $\rho$, $l\leq TP(\rho^{\prime})\leq U$, where $TP$ is simply the sum of the weight of the edges as defined earlier.
 
\chapter{Conclusion}
In this thesis, we have considered several variants of energy games. The first variant defines games with upper and lower bound constraints, combined with reachability or infinite runs objectives. 
The second variant proposed defines games with a strong lower bound and a relaxed upper bound that can be weak or temporarily exceeded, combined with reachability or infinite run objectives, and constraints on violations of upper bound. 
\vskip 0.3cm
In the one player case, complexities ranges from PTIME to PSPACE-Complete. and in the two-player case from NP $\cap$ coNP to EXPTIME-Complete. 
In general, the complexity is the same for a reachability and for an infinite run objective. Interestingly, for \LWenergy games, the complexity of the single player case is PTIME, but reachability objectives require exponential memory (in the size of the weak upper bound) while strategies are memoryless for infinite run objectives.
\vskip 0.3cm
Table~\ref{table-results} summarizes known
results, and the results obtained in this paper (where \LVenergyall
gathers all three energy constraints with violations). We furthermore show that the minimization problem for \LVenergyall (reachability) games require algorithms that run in PSPACE in the one-player case, and in EXPTIME in the two-players case.


\begin{table}[ht]
  \centering
  \def\arraystretch{2}
  \def\comp{c.}
  \scalebox{0.7}{
  \begin{tabular}{|c||c|c||c|c|}
    \hline
    &\multicolumn{2}{c|}{Reachability}&\multicolumn{2}{c|}{Infinite runs}\\
    \hline
    & 1 player & 2 players & 1 player & 2 players
    \\\hline\hline
    \Lenergy & PTIME~(thm.~\ref{thm_reachability_games})&NP$\cap$ coNP~(thm.~\ref{thm_reachability_games}) 
    & PTIME~\cite{BouyerFLMS08} & in NP $\cap$ coNP~\cite{BouyerFLMS08}\\\hline
    %LUEnergy Games
    \LUenergy&PSPACE-\comp~(thm.~\ref{pspace-complete})& EXPTIME-\comp ~(thm.~\ref{exp-complete})
    & PSPACE-\comp~\cite{BouyerFLMS08} &EXPTIME-\comp~\cite{BouyerFLMS08}\\\hline
    %
    \LWenergy&PTIME~(thm.~\ref{one-player-weak-thm})&coNP~(thm.~\ref{two-player-weak-thm}) 
    &PTIME~\cite{BouyerFLMS08} &NP $\cap$ coNP~\cite{BouyerFLMS08}\\\hline
    %
    APNA$^{*}$-games &PSPACE-\comp~(thm.~\ref{thm_apna_LUHV_2P}) &EXPTIME-\comp~(thm.~\ref{thm_apna_LUHV_2P})&PSPACE-\comp~(thm.~\ref{thm_apna_LUHV_2P}) &EXPTIME-\comp~(thm.~\ref{thm_apna_LUHV_2P}) \\\hline
    
    
    Bound existence&PSPACE-\comp~(thm.~\ref{thm_Apnar_exists_min_exptimec})&EXPTIME-\comp~(thm.~\ref{thm_Apnar_exists_min_exptimec}) &PSPACE-\comp~(thm.~\ref{thm_Apnar_exists_min_exptimec})&EXPTIME-\comp~(thm.~\ref{thm_Apnar_exists_min_exptimec}) \\\hline
  \end{tabular}
  }
  \caption{Summary of the results}
  \label{table-results}
\end{table}
\vskip 1cm
A possible extension of this work is to consider energy games with mean payoff function and discounted total payoff for the energy level and for the violation constraints, and the associated minimization and existence problems. 




\bibliographystyle{alpha}
\bibliography{Thesis}

\appendix
\chapter{Appendix Title}
\section{Proofs of Section~\ref{sec-weak}}

\noindent{\bf Lemma \ref{lemma-higherrun}: }
  Let $\pi$ be a finite path in a one-player arena~$G$.
  If $(q,u) \xrightarrow{\pi}_{LW} (q',u')$,
  then for any~$v\geq u$,
  $(q,v) \xrightarrow{\pi}_{LW} (q',v')$ for some~$v'\geq u'$.

\begin{proof}
Write $\pi=(e_i)_{0\leq i<n}$, with $e_i=(q_i,p_i,q'_i)$ for each~$i$.
The sequence defined as
\begin{xalignat*}2
  u_0&= u & u_{i+1} &= \min{W,u_i+p_i}
\end{xalignat*}
is the sequence of energy levels along the run $(q,u) \xrightarrow{\pi}_{LW}
(q',u')$.  For~$v\geq u$, letting
\begin{xalignat*}2
  v_0&= v & v_{i+1} &= \min{W,v_i+p_i},
\end{xalignat*}
we~easily prove by induction that for all~$i$, $u_i\leq v_i\leq W$,
which entails that $(q,v) \xrightarrow{\pi}_{LW} (q',v')$ with $v'=v_n\geq
u_n=u'$.
\end{proof}


\noindent{\bf Proof of Lemma~\ref{lemma-hitW}}
  Let $\pi$ be a finite path in a one-player arena~$G$, and consider two LW-runs 
  $(q,u)\xrightarrow{\pi}_{LW} (q',u')$ and $(q,v)\xrightarrow{\pi}_{LW}(q',v')$ with $u\leq
  v$. Then $u'-u\geq v'-v$, and if the inequality is strict, then
  the energy level along the run $(q,v)\xrightarrow{\pi}_{LW}(q',v')$ must have hit~$W$.


\begin{proof}
  The first statement is proven by induction: we~again write
  $\pi=(e_i)_{0\leq i<n}$, with $e_i=(q_i,p_i,q'_i)$ for each~$i$, and
  \begin{xalignat*}2
    u_0&= u & u_{i+1} &= \min(W,u_i+p_i)\\
    v_0&= v & v_{i+1} &= \min(W,v_i+p_i).
  \end{xalignat*}
  Then $u_{i+1}-u_i=\min(W-u_i,p_i)$ and
  $v_{i+1}-v_i=\min(W-v_i,p_i)$.  Since $u_i\leq v_i$ for all~$i$,
  we~also have $W-u_i\geq W-v_i$, and $u_{i+1}-u_i \geq
  v_{i+1}-v_i$. By~summing up these inequalities, we~get $u_{i+1}-u_0
  \geq v_{i+1}-v_0$. Now, as long as $W-v_i\geq p_i$ (then also
  $W-u_i\geq p_i$), the inequalities above are equalities. It~follows
  that if the inequality is strict, then the run
  $(q,v)\xrightarrow{\pi}_{LW}(q',v')$ must have hit~$W$.
\end{proof}


\noindent{\bf Proof of Lemma~\ref{lemma-W}}
Let $\pi$ be a finite path in a one-player arena~$G$, for which there
  is an LW-runs $(q,u)\xrightarrow{\pi}_{LW} (q',u')$. 
  %with $u\leq u'$. (suppressed this info otherwise the second statement makes no sense) 
  If $u'$ is the maximal energy level along that run, then
  $(q,W)\xrightarrow{\pi}_{LW} (q',W)$;
  if $u$ is the maximal energy level along the run above, then
  $(q,W)\xrightarrow{\pi}_{LW} (q',W+u'-u)$.


\begin{proof}
  Write $\pi=(e_i)_{0\leq i<n}$, with $e_i=(q_i,p_i,q'_i)$ for
  each~$i$. If~$u'$ is the maximal energy level, then for all~$i$,
  we~have $\sum_{j=i}^{n-1} p_j\geq 0$.  Now, define
  \begin{xalignat*}2
    v_0&= W & v_{i+1} &= \min(W,v_i+p_i).
  \end{xalignat*}
  If~$v_n<W$, then by induction we also have $v_i<W$ for all~$i$,
  contradicting the fact that~$v_0=W$. This proves our first result.

  Similarly, if~$u$ is the maximal energy level, then for all~$i$,
  we~have $\sum_{j=0}^{i} p_j\leq 0$. Then for all~$i$,
  $v_{i+1}=v_i+p_i\leq W$, so that $v_{n}-v_0=u'-u$. Our second result
  follows.
\end{proof}


\noindent{\bf Proof of Lemma~\ref{lemma-iteratepos}}
Let $\pi$ be a finite path in a one-player arena~$G$.
  If
  $(q,u)\xrightarrow{\pi}_{LW} (q',u')$ with $u'> u$
  and 
  $(q,w)\xrightarrow{\pi}_{LW} (q',w')$ with $w'> w$,
  then
  for any $u\leq v\leq w$, it~holds
  $(q,v)\xrightarrow{\pi}_{LW} (q',v')$ with $v'> v$.

\begin{proof}  
Using Lemma~\ref{lemma-higherrun}, we~immediately have
$(q,v)\xrightarrow{\pi}_{LW} (q',v')$. As~in the previous proof, we~define
sequences
\begin{xalignat*}2
  u_0&= u & u_{i+1} &= \min(W,u_i+p_i)\\
  v_0&= v & v_{i+1} &= \min(W,v_i+p_i)\\
  w_0&= w & w_{i+1} &= \min(W,w_i+p_i).
\end{xalignat*}
We~still have $u_i\leq v_i\leq w_i$ for all~$i$. Moreover, 
%if $v_j<M$
if $v_j<\wub$
for all~$j\leq i$, then $v_i-u_i=v-u$. As~a consequence, if $v'\leq
v$, then it must be the case that 
%$v_j=M$ 
$v_j=\wub$ 
for some~$j$; but then $w_j=v_j$, 
%since $v_j\leq w_j\leq M$. 
since $v_j\leq w_j\leq \wub$. 
It~follows that $w_k=v_k$ for
all~$k\geq j$, so at the end of $\pi$ we have $w'=v'$. Assuming $v'\leq v$ raises a contradiction since we have $v'=w'>w\geq v$.
%This is a contradiction, since we assumed $v'\leq v$ and have $w'>w$ by hypothesis. 
Hence $v'>v$.
\end{proof}

\noindent{\bf Proof of Lemma~\ref{lemma-removeneg}}
Let $\pi$ be a finite path in a one-player arena~$G$.
  If $(q,u) \xrightarrow{\pi}_{LW} (q',u')$ and $\pi$ can be decomposed as
  $\pi_1\cdot\pi_2\cdot \pi_3$ in such a way that
  $(q,u) \xrightarrow{\pi_1}_{LW}
%  (s,v)\LWarrow{\pi_2}(s',v') \LWarrow{\pi_3}(q',u')$ with $v'\leq
  (s,v)\xrightarrow{\pi_2}_{LW}(s,v') \xrightarrow{\pi_3}_{LW}(q',u')$ with $v'\leq
  v$, then $(q,u)\xrightarrow{\pi_1\cdot\pi_3}_{LW} (q',u'')$ with $u''\geq
  u'$.

\begin{proof}
%Since $(s',v') \LWarrow{\pi_3}(q',u')$ and $v'\leq v$, by
Since $(s,v') \xrightarrow{\pi_3}_{LW}(q',u')$ and $v'\leq v$, by
Lemma~\ref{lemma-higherrun} 
%we also have $(s',v)
we also have $(s,v)
\xrightarrow{\pi_3}_{LW}(q',u'')$ for some~$u''\geq u'$. The~result follows.
\end{proof}


\noindent{\bf Proof of Lemma~\ref{lemma-iteratecycles}}
Let $\pi$ be a cycle on~$q$ such that $(q,u) \xrightarrow{\pi}_{LW} (q,v)$
  for some $u\leq v$. Then $(q,u) \xrightarrow{\pi^{W-L}}_{LW} (q,v')$ for
  some~$v'$, and $(q,v')\xrightarrow{\pi}_{LW} (q,v')$.


\begin{proof}
  The case where $u=v$ is trivial. We~assume $u<v$.  Applying
  Lemma~\ref{lemma-higherrun} inductively, we get that the cycle can
  be iterated arbitrarily many times; this~also proves that the
  sequence of energy levels reached at the end of each iteration is
  non-decreasing.

  Now, assume that $(q,v')\xrightarrow{\pi}_{LW} (q,v'')$ for
  some~$v''\not=v'$.  Then $v''>v'$. Lemma~\ref{lemma-iteratepos} then
  entails that the sequence of energy levels reached at the end of
  each iteration is increasing. Since the loop has been iterated $W-L$
  times, the~energy level in~$v''$ would exceed~$W$, which is
  impossible. This proves our result.
\end{proof}


\noindent{\bf Proof of Lemma~\ref{lemma-DAGlabel}}
  Let~$[q,d]$ be a state of the DAG, and $M$ and~$m$ be two integers
  such that $0\leq m\leq M$.  Upon termination of this algorithm, 
  state~$[q,d]$ of the DAG is labelled with~$(M,m)$ if, and only if,
  there is an LW-run of length~$d$ from~$(q_0,L)$ to~$(q,M-m)$ along
  which the energy level always remains in the interval~$[L,M]$.

\begin{proof}
  The proof is by induction on~$d$. The result is trivial for~$d=0$.
  Now, assume it holds for some depth~$d-1$, and pick a
  state~$[q,d]$. For the first direction, if~$[q,d]$ is labelled
  with~$(M,m)$, then this label was added using some
  transition~$([q',d-1],w,[q,d])$ and some label~$(M',m')$
  of~$[q',d-1]$. By~induction, there is an LW-run~$\rho$ of
  length~$d-1$ from~$(q_0,L)$ to~$(q',M'-m')$ in~$G$ along which the
  energy level remains in the interval~$[L,M']$. We~consider two cases,
  corresponding to the two ways of updating the pair of values:
  \begin{itemize}
  \item if $w>m'$, then we~have $M=min(\wub,M'-m'+w)$ and $m=0$. %$M=\max(W,M'-m'+w)$ and $m=0$. 
  Now, the transition $([q',d-1],w,[q,d])$ in the DAG originates from a
    transition~$(q',w,q)$ in~$G$; taking this transition after~$\rho$
    provides us with the run of length~$d$ from~$(q_0,L)$ to~$(q,M-m)$
    along which the energy level remains in~$[L,M]$, as required;
  \item if $m'+L-M'\leq w\leq m'$, then $M=M'$ and $m=m'-w$. Again,
    taking transition~$(q',w,q)$ after~$\rho$ provides us with the
    LW-run we are looking for.
  \end{itemize}

  Conversely, if there is an LW-run~$\rho$ of length~$d$
  from~$(q_0,L)$ to~$(q,M-m)$ along which the energy level always
  remains in the interval~$[L,M]$, then we~write $\rho=\rho'\cdot
  ((q',l'),w,(q,M-m))$, distinguishing its last transition. By~induction, $[q',d-1]$ must have been labelled with a pair~$(M',m')$ such that $l'=M'-m'$ and the energy level
  along $\rho'$ remained within $[L,M']$. Now, from the existence of a transition $((q',l'),w,(q,M-m))$, we~know that there is a
  transition~$([q',d-1],w,[q,d])$ in the~DAG, which will generate the required label of~$[q,d]$.
\end{proof}

\noindent{\bf Proof of Lemma~\ref{lemma-univcycle}}
  Let~$[q_0,d]$ be a state of the DAG, with~$d>0$. Let $m$ be a
  non-negative integer such that $L<W-m$.  Upon termination of this algorithm, state~$[q_0,d]$ is labelled with~$(M,m)$ such that $M-m>L$ if, and only if, there is a universal  positive cycle~$\phi$ on~$q_0$ of length~$d$ such that $(q_0,L) \xrightarrow{\phi^{W-L}}_{LW} (q_0,W-m)$.

\begin{proof}
First assume that~$[q_0,d]$ is labelled with~$(M,m)$ for some~$M$ such
that $M-m>L$. From Lemma~\ref{lemma-DAGlabel}, there is a cycle~$\phi$
on~$q_0$ of length~$d$ generating a run $(q_0,L)\xrightarrow{\phi}_{LW}(q_0,M-m)$
along which the energy level is within~$[L,M]$.  Then~$M-m\geq L$, so
that Lemma~\ref{lemma-iteratecycles} applies: we~then get $(q_0,L)
\xrightarrow{\phi^{W-L}}_{LW} (q_0,E)$ with $(q_0,E)\xrightarrow{\phi}_{LW}(q_0,E)$.
Write $(p_i)_{0\leq i<|\phi|}$ for the sequence of weights along~$\phi$.
Also write $\rho$ for the run $(q_0,L)\xrightarrow{\phi}_{LW}(q_0,M-m)$, and
$\sigma$ for the run $(q_0,E)\xrightarrow{\phi}_{LW}(q_0,E)$.

%Assume $L=M-m$: iterating~$\phi$ would not modify the energy level
%reached at the end of those cycles, so that we would have
%$L=W-m$. Since this is not the case, it~must be 
As $L<M-m$, then by Lemma~\ref{lemma-hitW}, it~must be the case that energy level~$W$ is
reached along~$\sigma$.  Write~$i_0$ and $j_0$ for the first and last
positions along~$\rho$ for which the energy level along~$\rho$ is~$M$.
That is, the subpath from index $0$ to index $i_0$ has growing energy level, and the subpath from $j_0$ to $|\rho|$ has a decreasing energy level. 
Assume $\tilde\sigma_{i_0}\not=W$: by~Lemma~\ref{lemma-higherrun},
we~must have $M=\tilde\rho_{i_0}\leq \tilde\sigma_{i_0}<W$.  Then for
all~$k\geq i_0$, $\sum_{l=i_0}^k p_l\leq 0$, and $\sum_{l=i_0}^{j_0}
p_l= 0$. Since $\tilde\sigma_{i_0}<W$, then also $\tilde\sigma_{k}<W$
for all~$k\geq i_0$. According to Lemma~\ref{lemma-hitW}, energy level $\wub$ is reached in $\sigma$, so 
%HERE 
there exists some~$k_0<i_0$ such that $\tilde\sigma_{k_0}=W$;
%$\tilde\rho_{k_0}<M$, 
However, as $i_0$ is the index of the first maximal value in $\rho$, we have $\tilde\rho_{k_0}<M$, and the energy level increases in run $\rho$ between $k_0$ and $i_0$. So according to Lemma~\ref{lemma-W}, we should have $\tilde\sigma_{i_0}<W$, which raises a contradiction.
Hence we proved $\tilde\sigma_{i_0}=W$; applying
the second result of Lemma~\ref{lemma-W}, we~get $E=W-m$.

\medskip
Conversely, if there is a universal positive cycle~$\phi$ satisfying the
conditions of the lemma, let $F$ be such that $(q_0,L)\xrightarrow{\phi}_{LW}
(q_0,F)$, and~$M$ be the the maximal energy level encountered along
the run $(q_0,L)\xrightarrow{\phi}_{LW} (q_0,F)$. By
Lemma~\ref{lemma-DAGlabel}, state~$[q_0,d]$ is labelled with~$(M,m')$
for some~$m'\geq 0$ such that $F=M-m'$.
By~Lemma~\ref{lemma-iteratecycles}, we~must have $(q_0,L)
\xrightarrow{\phi^{W-L}}_{LW} (q_0,W-m')$.
%, so that $m'=m$.
\end{proof}


\noindent{\bf Proof of Lemma~\ref{lemma-order-cycle}}
Consider two paths~$\pi$ and~$\pi'$ such that
$first(\pi)=first(\pi')$ and $last(\pi)=last(\pi')$, and with
respective values~$(M,m)$ and~$(M',m')$ such that $(M,m)\preceq (M',m')$.
If~$\pi$ is a prefix of a universal cycle~$\phi$, then $\pi'$~is a
prefix of a universal cycle~$\phi'$ with $\phi'\triangleright\phi$.

\begin{proof}
  Let $q=first(\pi)$ and $q'=last(\pi)$.  We~write~$\psi$ for the
  path such that $\phi=\pi\cdot\psi$; $\psi$~is a path from~$q'$
  to~$q$. Then $(q,L)\xrightarrow{\pi}_{LW} (q'',M-m)
  \xrightarrow{\psi}_{LW}(q',F)$. Also, $(q,L)\xrightarrow{\pi}_{LW} (q'',M'-m')$.  Since
  $M-m\leq M'-m'$, we have $(q'',M'-m')\xrightarrow{\psi}_{LW}(q',F')$. We~can
  thus let $\phi'=\pi'\cdot\psi$: by~Lemma~\ref{lemma-univcycle},
  the~final energy level reached after iterating~$\phi'$ is higher
  than the energy level reached after iterating~$\phi$, since $m'\leq
  m$. Hence $\phi'\triangleright\phi$.
\end{proof}



\noindent{\bf Proof of Lemma~\ref{lemma-DAG-construction-poly}}
%If in our algorithm we only store maximal labels (for~$\preceq$), then
%the algorithm runs in polynomial time.
If the DAG construction algorithm only stores maximal labels (for~$\preceq$), then
it runs in polynomial time.

\begin{proof}
We prove that, when attaching to each node $[q,d]$ of the DAG only the maximal labels (w.r.t $\preceq$) reached for a path of length $d$ ending in state $q$, the number of
values for the first component of the different labels that appear at
depth~$d>0$ in the DAG is at most $d\cdot |Q|$. Since it only
stores optimal labels, our algorithm will never associate to a state $[q,d]$
two labels having the same value on their first component. So, any
state at depth~$d$ will have at most $d\cdot |Q|$ labels.

So we prove, by induction on~$d$, that the number of different values
for the first component among the labels appearing at depth~$d>0$ is
at most~$d\cdot|Q|$. 
%
This is true for~$d=1$ since the initial state~$(q,0)$ only
contains~$(M=0,m=0)$, and each transition with nonnegative weight~$w$
will create one new label~$(w,0)$ (transitions with negative weight
are not prefixes of universal cycles). Now, since all those labels
%have the same value on their second component,
have value 0 as second component, each state $[q,1]$ in the DAG will be attached at most one label. Hence, the total number of labels (and the total number of different values for their first
component) is at most~$|Q|$ at depth 1 in the DAG.

Now, assume that the labels appearing 
%at depth~$q$ 
at depth $d>1$ are all drawn from a set of labels~$L=\{M_i,m_i) \mid 1\leq i\leq n\}$ in which the
number of different values of~$M_i$ is at most $d\cdot|Q|$.
Consider a state~$[q',d]$, labelled with $\{(M_i,m_i) \mid 1\leq i\leq
n_{q',d}\}$ (even if it means reindexing the labels). Pick~a
transition from~$[q',d]$ to~$[q'',d+1]$, with weight~$w$. For each
pair~$(M_i,m_i)$ associated with $[q,d]$, it~creates a new label in~$(q'',d+1)$; this label
is
\begin{itemize}
\item either $(M_i-m_i+w,0)$ if~$m_i<w$;
\item or $(M_i,m_i+w)$ if $m_i=M_i\leq w\leq m_i$.
\end{itemize}
Now, for a state~$(q'',d+1)$, the~set of labels created by all
incoming transitions can be grouped as follows:
\begin{itemize}
\item labels having~zero as their second component; among those, our
  algorithm only stores the one with maximal first component, as $(M_i,0) \preceq (M_j,0)$ as soon as $M_i\leq M_j$;
\item for each $M_i$ appearing at depth~$d$, labels having~$M_i$ as
  their first component; again, we only keep the one with minimal
  second component, as $(M,m_i) \preceq (M,m_j)$ when $m_j \leq m_i$.
\end{itemize}
In~the end, for this state~$[q'',d+1]$, we keep at most one label for
each distinct value among the first components $M_i$ of labels appearing at depth~$d$, and possibly one extra
label with second value~$0$. In~other terms, at depth $d+1$ the values that appear as first component of labels are 
obtained from values at depth~$d$, plus possibly one per state; Hence, at depth $d+1$, there exists at most $(d+1)\cdot|Q|$ labels, which completes the proof of the induction step.
\end{proof}



\noindent{\bf Proof of Lemma~\ref{lemma-higherstrat}}
  Let $G$ be a two-player arena, equiped with an \LWenergy-reachability
  objective. Let~$q$ be a state of~$G$, and~$u\leq u'$ in $[L;W]$. If
  \Pl1 wins the game from~$(q,u)$, then she also wins from~$(q,u')$.

\begin{proof}
Let~$\sigma$ be a winning strategy for \Pl1 from~$(q,u)$. If she plays
the same strategy from~$(q,u')$, then for any strategy of~\Pl2, the
resulting outcome from~$(q,u')$ follows the same transitions as the
outcome of the same strategies from~$u$, with higher energy
level. Since~$\sigma$ is winning from~$(q,u)$, it~is also winning
from~$(q,u')$.
\end{proof}

\end{document}
